\chapter{持久化、崩溃一致性与恢复}

\section{易失状态与持久状态}

持久化问题的核心是：系统可能在任意两步之间崩溃。程序调用 \code{write}，内核可能只把数据放进 page cache；设备可能只完成了一部分请求；目录项和文件内容可能以不同顺序持久化；应用日志可能写了意图但没有应用数据；manifest 可能发布了一个并不完整的输出。崩溃一致性讨论的不是“正常运行时看起来对不对”，而是“重启后能否解释磁盘状态”。

\code{write}、\code{fsync}、\code{rename} 和 manifest 分别处在不同层次。\code{write} 把字节交给内核；\code{fsync} 推进文件内容和部分元数据的持久化；\code{rename} 可以提供目录项替换的原子可见性，但目录本身是否持久也需要考虑；manifest 是业务层承认某批输出有效的提交点。只有把这些层次分开，才能写可靠输出。

文件系统常用日志或 copy-on-write 技术维护自己的元数据一致性，但这不自动等于应用数据一致。文件系统保证的是它自己的结构能恢复，不保证你的两个文件之间满足业务事务。应用如果需要“数据文件和索引文件同时生效”，就必须设计应用层提交协议。

最小可靠写出协议通常是：写临时文件，写完后 \code{fsync} 临时文件，校验大小和 checksum，\code{rename} 到最终路径，\code{fsync} 父目录，最后写或更新 manifest。这个协议不是唯一答案，但它展示了 prepared 和 committed 的区别。

日志恢复要区分 redo 和 undo。redo 日志记录“应该完成什么”，崩溃后可以重放；undo 日志记录“如何撤销未完成修改”。write-ahead logging 的关键规则是：描述修改的日志必须先于实际修改持久化，否则崩溃后无法判断磁盘上的半成品是否应该保留。

\section{从一次写入逐步得到持久化协议}

先不使用同步、日志和事务这些名称。设程序要把文件中的整数从 41 改为 42。最直接的
实现是把新字节交给 \code{write}，看到成功返回后就向用户报告完成。这个结论只有在
系统一直运行时才可能成立；如果返回之后立即掉电，重启后仍可能读到 41。

原因在于一次写入会经过多个保存位置。用户缓冲区中的 42 首先被复制或引用，随后进入
内核缓存；内核可以稍后才产生设备请求；设备也可能先把请求放进自己的易失缓存，最后才
使介质上的内容具备跨掉电保存的能力。每一层都可能说“我已经接收”，但“接收”与
“以后仍然存在”不是同一个承诺。

\begin{lstlisting}
user buffer
  -> kernel page cache
  -> filesystem writeback request
  -> device queue or volatile cache
  -> non-volatile media
\end{lstlisting}

因此，讨论“写成功”时至少要回答四个不同问题：

\begin{longtable}{@{}F{0.19\linewidth}F{0.36\linewidth}F{0.35\linewidth}@{}}
\toprule
\rowcolor{BookBlueTint}
问题 & 所要求的事实 & 不能据此推断的结论 \\
\midrule
可见性 & 当前或其他执行流已经能读到新值 & 掉电后仍能读到新值 \\\tablerowrule
顺序性 & 更新 A 在更新 B 之前达到指定边界 & A 与 B 作为整体原子生效 \\\tablerowrule
持久性 & 已确认的更新能在规定故障模型下保留 & 多个更新之间保持业务一致 \\\tablerowrule
原子性 & 观察者只看到操作前或操作后的完整状态 & 新状态一定已经持久化 \\
\bottomrule
\end{longtable}

这四项不能互相替代。\code{rename} 可以使路径替换具有原子可见性，却不自动证明目录
修改已经跨过持久化边界；\code{fsync} 可以推动一个文件的内容和相关元数据持久化，
却不能自动把两个独立文件合成一项业务事务；日志可以使文件系统结构恢复到一致状态，
却不知道应用所说的“版本 42”还包含哪些文件。

\fstopic{为什么需要显式的持久化边界}

如果所有写都立即等待介质完成，语义会简单一些，但每一次很小的修改都要承担设备延迟，
顺序写合并和缓存复用也难以发挥作用。操作系统因此允许普通写先在内存中完成，把昂贵的
设备提交推迟或合并。性能由此提高，同时也产生一个必须显式表达的边界：应用何时确实需要
等待数据跨过易失层。

对普通文件来说，这个边界通常由 \code{fsync} 或更具体的同步接口表达。它不是“再写一次”，
而是要求文件系统找出与该文件有关的脏数据和必要元数据，发出写回，等待错误可判定地返回，
并在设备支持范围内传播刷新或强制单元访问要求。任意一层忽略这项要求，上层看到的成功就会
失去原有含义。

这里还要区分文件内容和名称关系。文件同步主要围绕 inode、数据块和必要元数据；新建、
删除或改名改变的是父目录保存的名字映射。若程序希望“重启后这个名字一定存在并指向新文件”，
只同步文件内容可能不够，还要把目录项变化推进到持久化边界。这就是临时文件发布协议在
\code{rename} 之后继续同步父目录的原因。

\fstopic{为什么写入顺序也需要证据}

设应用先写数据块 D，再写一个指向 D 的元数据指针 M。程序发出 D 后再发出 M，并不能保证
介质按同样顺序完成：内核、块层、控制器和设备都可能为了吞吐重新排序。如果 M 先持久化，
随后掉电，重启后元数据会指向尚未写好的 D。

正确协议必须使“D 已达到要求的边界”成为发布 M 的前提。可以等待 D 完成并刷新后再写 M；
也可以先把这次修改的描述写入日志并提交；还可以用写时复制把新数据和新索引写到未发布位置，
最后切换可恢复的根引用。三种方法的结构不同，但都在增加同一种东西：恢复程序能够依赖的
顺序证据。

\begin{lstlisting}
unsafe:
  issue(D)
  issue(M points to D)
  // issue order alone is not durable order

ordered publication:
  write(D)
  wait_until_D_is_durable()
  write(M points to D)
  persist(M)
\end{lstlisting}

\fstopic{为什么多对象更新需要恢复协议}

单个文件创建已经会改变 inode 分配状态、inode 内容、目录项和可能的数据块分配状态。
这些状态不可能依靠一次普通内存赋值同时写入介质。若系统在第二项和第三项之间崩溃，介质
就会保存一个正常运行时不应长期出现的中间状态。

系统有两种基本选择。第一种是在重启后扫描全局结构，根据不变量判断并修复中间状态；
第二种是在正常更新时额外保存意图与提交证据，使恢复只处理最近尚未归并的修改。前者得到
全盘一致性检查，恢复工作量通常与文件系统规模有关；后者逐步得到日志、事务与检查点，
恢复工作量主要与尚未归并的记录有关。

这时才有必要引入\emph{事务}一词。这里的事务不是数据库功能的简单借用，而是一组必须以
同一个恢复结论处理的更新。提交记录表示这组更新已经取得发布资格；没有提交记录的尾部
不能被恢复程序当作已完成事实。若采用 redo，已提交但尚未全部写回的更新要能够重复执行；
若采用 undo，未提交却已经泄漏到目标位置的修改要能够撤销。

\fstopic{故障模型决定保证范围}

“绝不丢数据”不是可以脱离条件成立的保证。系统必须先规定考虑哪些故障：进程退出、内核
崩溃、整机掉电、控制器丢失缓存、介质返回错误、单盘永久损坏，还是多个副本同时损坏。
日志可以处理更新中断，却不能从唯一一块已经永久损坏的介质中恢复原数据；镜像可以承受
一个副本丢失，却不能自动判断两个副本中哪个包含更新后的正确字节；校验和能够识别内容与
期望不符，却必须配合冗余才有机会修复。

因此，后文所有恢复算法都要同时写明三件事：它把哪些记录视为权威事实，允许介质停留在
哪些中间状态，以及超出故障模型后选择拒绝挂载、只读降级、丢弃事务还是请求人工处理。

\section{文件发布、目录同步与 manifest}

先把一次发布过程写成提交状态机：

\begin{lstlisting}
NoOutput
  -> TempCreated
  -> DataWritten
  -> DataSynced
  -> RenamedVisible
  -> DirectorySynced
  -> ManifestCommitted
\end{lstlisting}

每个状态都要写出崩溃后恢复规则：

\begin{longtable}{p{0.24\linewidth}p{0.66\linewidth}}
\toprule
崩溃点 & 恢复策略 \\
\midrule
TempCreated & 删除临时文件或忽略 \\
DataWritten & 校验失败则删除，不能发布 \\
DataSynced & 仍未提交，等待 rename 或清理 \\
RenamedVisible & 若 manifest 未提交，按业务规则决定是否采用 \\
DirectorySynced & 文件名持久，但业务仍以 manifest 为准 \\
ManifestCommitted & 读取者可以信任并校验文件 \\
\bottomrule
\end{longtable}

实践中要避免把“文件存在”当成业务成功。目录中有最终文件名，只说明某次 rename 可能发生过；manifest 提交并通过 checksum，才说明业务层承认它属于某个版本或批次。

再写一份日志协议：

\begin{lstlisting}
append log record: transaction T intends to write block B with checksum C
fsync log
write data block
fsync data
append commit record for T
fsync log
\end{lstlisting}

恢复时扫描日志：有 commit 的事务必须 redo；没有 commit 的事务按规则忽略或 undo；checksum 不匹配的记录停止或进入人工恢复。这里的重点不是实现数据库，而是理解日志先行、提交记录和幂等恢复。

\section{WAL、检查点与幂等恢复}

\subsection{原子 rename 发布}

很多应用用临时文件加 rename 发布结果。算法分为写入、同步、替换、同步目录四步：

\begin{lstlisting}
write_atomic(path, bytes):
  tmp = path + ".tmp." + unique_id()
  fd = open(tmp, O_CREAT | O_EXCL | O_WRONLY)
  write_all(fd, bytes)
  fsync(fd)
  close(fd)
  rename(tmp, path)
  dirfd = open(parent_dir(path))
  fsync(dirfd)
  close(dirfd)
\end{lstlisting}

这个算法保证读者要么看到旧文件，要么看到新文件，不应看到半个最终文件。但它仍然只是单文件发布。若业务需要多个文件同时生效，还需要 manifest 或事务日志。

\subsection{manifest 提交}

manifest 把多个输出文件绑定成一个版本。数据文件可以先各自写好并 fsync；manifest 最后写入，列出版本号、文件名、大小和 checksum。读取者只读取 manifest 指向的文件。

\begin{lstlisting}
commit_version(version, files):
  for file in files:
    write_atomic(file.tmp_path, file.bytes)
    verify_checksum(file.tmp_path)
    rename(file.tmp_path, file.final_path)
  manifest = build_manifest(version, files)
  write_atomic("MANIFEST.new", manifest)
  rename("MANIFEST.new", "MANIFEST")
\end{lstlisting}

manifest 的优点是恢复简单：重启后读取 MANIFEST，把它当作权威事实；不在 manifest 里的临时文件或孤儿文件可以清理。缺点是提交粒度较粗，manifest 本身也需要可靠写出。

\subsection{redo 日志}

redo 日志适合“提交后必须能重放”的场景。日志记录足够的信息，使恢复程序可以重新执行已经提交但未完全应用的修改。日志记录必须先持久化，数据页才能写出。

\begin{lstlisting}
transaction_write(T, block, data):
  log_append(INTENT, T, block, checksum(data))
  fsync(log)
  write(block, data)
  log_append(COMMIT, T)
  fsync(log)
\end{lstlisting}

恢复算法扫描日志：看到 INTENT 但没有 COMMIT 的事务不应用；看到 COMMIT 的事务检查数据是否已经应用，未应用则 redo。redo 操作必须幂等，即重复执行不会改变最终正确结果。

\subsection{检查点}

日志无限增长不可接受，所以系统要做 checkpoint。checkpoint 把当前稳定状态写成快照，并记录“恢复只需从这个日志位置之后开始”。算法难点是 checkpoint 期间仍可能有新事务进入，因此要记录 epoch 或 LSN。

\begin{lstlisting}
checkpoint():
  begin_lsn = log.current_lsn()
  flush_dirty_pages_up_to(begin_lsn)
  write_checkpoint_record(begin_lsn)
  fsync(log)
  truncate_log_before(begin_lsn)
\end{lstlisting}

checkpoint 的正确性依赖顺序：只有当某个 LSN 之前的脏数据都已经安全写出，才能截断之前的日志。否则崩溃后会丢失 redo 所需信息。

\section{撕裂写、校验和与介质故障}

持久化测试必须做故障注入。只在正常路径写文件再读取，不足以证明崩溃一致性。

\begin{itemize}
\tightlist
\item 在每个提交状态后模拟崩溃，恢复程序能给出确定结果。
\item 临时文件残留不会被读取者当成有效输出。
\item manifest 指向缺失文件时，读取者拒绝该版本。
\item checksum 错误时，恢复程序能定位坏文件并拒绝提交。
\item 重复执行恢复程序是幂等的，不会一次恢复成功、二次恢复破坏状态。
\item \code{write} 成功但 \code{fsync} 失败时，业务提交必须失败。
\item rename 后但 manifest 前崩溃，恢复策略必须明确。
\end{itemize}

测试报告不要只写“断电恢复成功”。必须列出崩溃点、磁盘上允许出现的文件集合、恢复后权威事实、被清理对象和未覆盖边界。可靠性来自状态枚举，不来自运气。

\section{COW、RAID 与多层存储映射}

\subsection{崩溃模型分层}

崩溃一致性至少有三类故障：transaction crash，系统在两步之间断电；system crash，内核或整机停止但介质仍可靠；media failure，介质本身丢失或返回坏数据。日志能处理前两类的一部分，但不能让坏盘自动恢复数据，后者需要冗余、校验和副本。

\begin{longtable}{p{0.2\linewidth}p{0.36\linewidth}p{0.32\linewidth}}
\toprule
故障 & 例子 & 需要的机制 \\
\midrule
transaction crash & commit 前断电 & WAL/journal/manifest 状态机 \\
system crash & 内核 panic、掉电 & fsync、flush、replay、fsck \\
media failure & 坏块、整盘损坏 & RAID、复制、checksum、备份 \\
\bottomrule
\end{longtable}

\subsection{LSN 与 checkpoint}

WAL 系统通常给每条日志记录 LSN。数据页记录 pageLSN，表示该页包含到哪个日志位置的修改。恢复时若 pageLSN 小于 committed LSN，就 redo；checkpoint 记录某个 LSN 之前的脏页已经安全传播。

\begin{lstlisting}
redo_recovery():
  start = last_checkpoint_lsn
  for record in log from start:
    if record.committed and page.page_lsn < record.lsn:
      apply(record)
      page.page_lsn = record.lsn
\end{lstlisting}

LSN 的价值是让恢复可判定，而不是靠文件时间戳或猜测。数据库、文件系统 journal 和 manifest 版本号都在解决同一问题：给恢复程序一个可比较的事实。

\subsection{COW 文件系统}

btrfs、ZFS 等 COW 文件系统不在原地覆盖元数据，而是写新块，再用新的根指针发布版本。snapshot、clone、send/receive 都建立在“旧版本块仍被引用，新版本写到新位置”的基础上。

\begin{lstlisting}
cow_update(tree, key, value):
  new_leaf = copy_and_modify(old_leaf)
  new_parent = copy_and_point_to(new_leaf)
  ...
  write_all_new_blocks()
  atomically_update_super_root(new_root)
\end{lstlisting}

COW 避免部分原地覆盖造成的结构破坏，但带来写放大、碎片和空间回收复杂度。snapshot 多时，一个旧块可能被许多版本引用，不能简单释放。

\subsection{RAID 与 write hole}

RAID 0 条带化提高吞吐但无冗余；RAID 1 镜像；RAID 5/6 用奇偶校验节省空间；RAID 10 结合镜像和条带。RAID 5 写入数据和 parity 时若断电，可能出现 write hole：数据和校验不一致。

\begin{longtable}{p{0.18\linewidth}p{0.36\linewidth}p{0.34\linewidth}}
\toprule
级别 & 优点 & 风险 \\
\midrule
RAID0 & 高吞吐、高容量利用 & 任一盘坏全丢 \\
RAID1 & 简单冗余，读可并行 & 容量减半 \\
RAID5/6 & 容量效率高，有校验 & 小写放大、write hole、重建压力 \\
RAID10 & 性能和可靠性较好 & 成本高 \\
\bottomrule
\end{longtable}

write hole 可用日志、非易失缓存、bitmap 或 copy-on-write parity 缓解。这里再次说明：持久化正确性来自顺序证据和冗余，不来自“写了两次应该差不多”。

\subsection{LVM 和 device mapper}

Linux device mapper 把上层块设备映射到底层设备区域，LVM 的 linear、snapshot、thin provisioning 都建立在此之上。dm-snapshot 用 COW 保存旧块，dm-thin 动态分配块，带来空间耗尽时的错误传播问题。

\begin{lstlisting}
bio to /dev/mapper/vg-lv
  -> dm target maps logical sector
  -> lower device sector(s)
  -> submit cloned bio
  -> complete original bio after all lower bios complete
\end{lstlisting}

设备映射层需要处理 bio 克隆、错误合并、flush 传播和快照元数据一致性。上层的持久化
请求只有被所有相关下层正确传播后，才能维持原有保证。
